\chapter{Data Understanding and Data Exploration}
\label{ch:data-understanding}

\section{Objectives of Data Understanding}

The data understanding phase (\textit{Data Understanding}) constitutes an indispensable prerequisite for any Data Mining activity. Before selecting or applying any algorithm, the analyst must acquire an overall view of the structure, distribution, and intrinsic quality of the dataset with respect to the project objectives.

In this preliminary phase, it is necessary to answer crucial methodological questions:
\begin{itemize}
    \item What types of attributes compose the dataset?
    \item What is the overall level of data quality (presence of noise, errors, or duplicates)?
    \item How are the numerical and categorical variables distributed?
    \item Are there evident correlations or redundancies among the attributes?
    \item How are outliers structured and what application meaning do they assume?
    \item What strategies should be adopted for missing values?
    \item Is it necessary to derive or transform attributes to make mining more effective?
\end{itemize}

\section{Types of Datasets}

The collected data present themselves in different logical-geometric structures based on the nature of the application domain:

\begin{enumerate}
    \item \textbf{Record Data:}
    \begin{itemize}
        \item \textbf{Data Matrix:} A set of objects described by the same fixed set of numerical attributes. A dataset of $m$ objects and $n$ attributes is formally represented by an $m \times n$ matrix:
        \begin{equation}
        \mathbf{X} = \begin{pmatrix}
        x_{1,1} & x_{1,2} & \cdots & x_{1,n} \\
        x_{2,1} & x_{2,2} & \cdots & x_{2,n} \\
        \vdots & \vdots & \ddots & \vdots \\
        x_{m,1} & x_{m,2} & \cdots & x_{m,n}
        \end{pmatrix}
        \end{equation}
        In this formalization, each object $x_i$ is a point in an $n$-dimensional vector space $\mathbb{R}^n$.
        \item \textbf{Document Data (Bag-of-Words):} Each textual document is represented by a term vector (\textit{term vector}). Each component of the vector corresponds to a word in the vocabulary, and the associated value indicates its frequency of occurrence in the text.
        \item \textbf{Transaction Data:} Each record represents a transaction (identified by a TID) and contains a variable set of purchased items or observed events.
    \end{itemize}

    \item \textbf{Graph Data:} Structures composed of nodes and edges that model relationships and connectivity. Examples include the network of hypertext links of the World Wide Web, social network graphs, and molecular structures (for example, the planar molecule of benzene $\text{C}_6\text{H}_6$).

    \item \textbf{Ordered Data:}
    \begin{itemize}
        \item \textbf{Sequences of transactions:} Temporal successions of purchases or events for each user.
        \item \textbf{Genomic data:} Discrete biological sequences composed of nucleotide characters $\{\text{A}, \text{C}, \text{G}, \text{T}\}$.
        \item \textbf{Spatio-Temporal Data:} Continuous or sampled measurements over geographical coordinates and regular time intervals (for example, the monthly average temperature of the land and ocean surface).
    \end{itemize}
\end{enumerate}

\section{Objects and Attributes}

A dataset is a collection of \textbf{data objects} and their corresponding \textbf{attributes}.
\begin{itemize}
    \item An \textbf{object} (also known as a record, point, case, sample, entity, or instance) represents a single entity of the domain of interest.
    \item An \textbf{attribute} (also known as a variable, field, characteristic, dimension, or feature) represents a measurable property or characteristic of an object (for example, age, taxable income, or eye color).
\end{itemize}

\subsection{Measurement Scales of Attributes}
The type of measurement scale of an attribute is defined by the set of algebraic properties and mathematical operators that preserve the meaning of the measurement:
\begin{enumerate}
    \item \textbf{Distinctness} ($=, \neq$): Possibility to distinguish whether two values are equal or different.
    \item \textbf{Order} ($<, >$): Possibility to establish an order relationship between values.
    \item \textbf{Meaningful differences} ($+, -$): Presence of a constant unit of measurement that allows quantifying the interval between two values (zero is conventional/arbitrary).
    \item \textbf{Meaningful ratios} ($*, /$): Presence of an \textbf{intrinsic absolute zero} (indicating the total absence of the measured quantity), which makes the proportional comparison between quantities meaningful.
\end{enumerate}

\begin{table}[htbp]
\centering
\footnotesize
\begin{tabularx}{\textwidth}{l p{2.2cm} X p{2.6cm} p{3.0cm}}
\toprule
\textbf{Scale} & \textbf{Properties} & \textbf{Description} & \textbf{Examples} & \textbf{Statistical Operators} \\
\midrule
\textbf{Nominal} & Distinctness ($=, \neq$) & Qualitative values or categorical labels without intrinsic ordering & SSN, ZIP code, eye color, gender & Mode, entropy, $\chi^2$ test, contingency \\
\addlinespace
\textbf{Ordinal} & Distinctness and Order ($=, \neq, <, >$) & Values with meaningful ordering but non-constant interval & Educational degree, Likert scale (1--10), mineral hardness & Median, percentiles, rank correlation (Spearman) \\
\addlinespace
\textbf{Interval} & Order and Differences ($+, -$) & Quantitative differences with meaning; zero is conventional & Temperature in $^\circ\text{C}$ or $^\circ\text{F}$, calendar dates & Mean, variance, standard deviation, Pearson, $t$ and $F$ tests \\
\addlinespace
\textbf{Ratio} & All 4 properties ($=, <, +, *, /$) & Presence of absolute zero; ratios between physically valid measurements & Age, income, mass, counts, Kelvin & Geometric and harmonic mean, percentage variation \\
\bottomrule
\end{tabularx}
\caption{Taxonomy of attribute measurement scales (Stevens).}
\label{tab:attribute-types}
\end{table}

\subsection{Further Distinctions of Attributes}
\begin{itemize}
    \item \textbf{Binary Attributes:} Nominal attributes that exclusively admit two states ($0$ and $1$). They are distinguished into:
    \begin{itemize}
        \item \textit{Symmetric Binary:} Both outcomes have equal informational importance (for example, biological gender).
        \item \textit{Asymmetric Binary:} One of the two outcomes is rare but of extreme informational relevance (for example, a positive outcome to a diagnostic test for a rare disease). By convention, the value $1$ is associated with the event of greater interest.
    \end{itemize}
    \item \textbf{Discrete vs Continuous Attributes:}
    \begin{itemize}
        \item \textit{Discrete:} Assume a finite or countably infinite set of values (typically represented by integers, counts, or categories). Binary attributes constitute a special case of discrete attribute.
        \item \textit{Continuous:} Assume values in a continuous subset of real numbers $\mathbb{R}$ (temperature, weight, coordinates). In computational practice, they are measured and stored using floating-point representations (\textit{floating-point}) with finite precision.
    \end{itemize}
\end{itemize}

\section{Data Quality}

Data quality is a critical factor for the success of any analytical initiative. As highlighted in Thomas Redman's studies, problems arising from poor quality data typically cost organizations between $10\%$ and $20\%$ of their overall revenues.
If a credit risk model is trained on flawed or erroneous data, the system will reject solvent debtors and disburse loans to insolvent individuals.

\subsection{Dimensions of Data Quality}
Quality problems manifest themselves under various dimensions:
\begin{enumerate}
    \item \textbf{Syntactic Accuracy:} The recorded value violates formal constraints or does not belong to the predetermined domain (for example, the string \texttt{"fmale"} for gender, or the presence of letters in a numeric field). Such anomalies can be identified immediately and automatically through simple schema and type checks.
    \item \textbf{Semantic Accuracy:} The data is formally correct and conforms to the expected type, but false in the reference reality (for example, associating the name \texttt{"John Smith"} with the female gender or recording a biological age of $180$ years). Detecting these discrepancies requires domain rules (\textit{business rules}) or cross-referencing with certified external sources.
    \item \textbf{Completeness:} Systematic lack of values for certain attributes or absence of entire records. An insidious violation of completeness is \textbf{sampling bias} (for example, the analysis of insolvency risk carried out on a historical archive that only includes customers already approved by the bank, excluding a priori rejected low-income profiles).
    \item \textbf{Unbalanced Data:} Strongly asymmetric distribution of classes (for example, fraudulent transactions equal to $0.1\%$ of the total or defective products on highly automated production lines).
    \item \textbf{Timeliness:} Obsolete degree of data updating compared to the current dynamics of the observed phenomenon.
    \item \textbf{Duplicate Data:} Presence of identical or nearly identical records relating to the same entity (for example, the same customer registered multiple times with different email addresses following the merger of two heterogeneous databases). Such issues require dedicated deduplication and \textit{record linkage} processes.
\end{enumerate}

\begin{noteblock}{Anomalies in Measurement Streams and Sensors}
In industrial monitoring and IoT systems, a broken or blocked sensor typically transmits a flat, constant value of zero. During the exploratory phase, such readings must not be treated as genuine physical measurements, but rather interpreted and handled in all respects as \textbf{missing values} (caused by hardware malfunction).
\end{noteblock}

\section{Descriptive Statistics and Univariate Exploration}

\subsection{Measures of Central Tendency}
\begin{itemize}
    \item \textbf{Sample Mean ($\bar{x}$):} For a sample of size $m$:
    \begin{equation}
    \bar{x} = \frac{1}{m} \sum_{i=1}^m x_i
    \end{equation}
    The mean is a distributive measure (it can be calculated by partitioning the dataset into subsets), but it is \textbf{extremely sensitive to the presence of outliers}.
    \item \textbf{Median:} The value that occupies the central position in the increasing order of the observations (or the mean of the two central values if $m$ is even). It is a robust measure, with a breakdown point (\textit{breakdown point}) of $50\%$.
    \item \textbf{Trimmed Mean:} Mean calculated by excluding a fixed percentage (for example $5\%$ or $10\%$) of the most extreme values at both ends.
    \item \textbf{Mode:} The value or category that occurs with the maximum absolute frequency. A distribution can be unimodal, bimodal, multimodal, or modeless if all observations appear only once.
\end{itemize}

\subsection{Skewness and Multimodality}
\begin{itemize}
    \item \textbf{Positively Skewed / Right-skewed:} The right tail of the distribution is elongated. Most observations are concentrated on values lower than the arithmetic mean ($\text{Mode} < \text{Median} < \text{Mean}$). A typical example is the distribution of \textbf{incomes}: the vast majority of citizens earn low or medium incomes, while a small number of individuals earn extremely high incomes, pulling the mean upwards.
    \item \textbf{Bimodal Distribution:} Presence of two distinct peaks in the frequency density. It typically signals the influence of an unstated binary latent variable (for example, the distribution of human body weight in a mixed population, influenced by the binary variable sex).
\end{itemize}

\subsection{Measures of Dispersion}
\begin{itemize}
    \item \textbf{Range:} Difference between maximum value and minimum value: $\text{Range} = x_{\max} - x_{\min}$.
    \item \textbf{Sample Variance ($s^2$) and Standard Deviation ($s$):}
    \begin{equation}
    s^2 = \frac{1}{m-1} \sum_{i=1}^m (x_i - \bar{x})^2, \quad s = \sqrt{s^2}
    \end{equation}
    The standard deviation measures the dispersion around the mean; it is zero only if all observations are identical ($s=0$), otherwise $s > 0$. It is also sensitive to outliers.
    \item \textbf{Robust Measures of Dispersion:}
    \begin{itemize}
        \item \textit{Absolute Average Deviation} (AAD):
        \begin{equation}
        \text{AAD} = \frac{1}{m} \sum_{i=1}^m |x_i - \bar{x}|
        \end{equation}
        \item \textit{Median Absolute Deviation} (MAD):
        \begin{equation}
        \text{MAD} = \operatorname{median}\big(|x_i - \operatorname{median}(x)|\big)
        \end{equation}
        \item \textit{Interquartile Range} (IQR): Distance between the third quartile $Q_3$ ($75^{\text{th}}$ percentile) and the first quartile $Q_1$ ($25^{\text{th}}$ percentile):
        \begin{equation}
        \text{IQR} = Q_3 - Q_1
        \end{equation}
    \end{itemize}
\end{itemize}

\subsection{Box Plot and Five-Number Summary}
The \textbf{Box Plot} (box and whisker plot) offers a non-parametric visual summary of the distribution based on the \textbf{five-number summary}:
\[
\big[\text{Minimum}, \; Q_1, \; \text{Median} \, (Q_2), \; Q_3, \; \text{Maximum}\big]
\]
Its geometric structure relies on the concept of quantile: given $p \in (0, 100)$, the \textbf{$p\%$-quantile} (or $p^{\text{th}}$ percentile) is the cutoff value $x_p$ such that $p\%$ of sample observations are smaller than or equal to $x_p$, while $(100-p)\%$ are greater. Specifically:
\begin{itemize}
    \item $Q_1$ corresponds to the $25\%$-quantile (first quartile).
    \item The median $Q_2$ corresponds to the $50\%$-quantile (second quartile).
    \item $Q_3$ corresponds to the $75\%$-quantile (third quartile).
    \item The width of the central box represents the interquartile range $\text{IQR} = Q_3 - Q_1$, encompassing the central $50\%$ of the data mass.
\end{itemize}

\begin{figure}[htbp]
\centering
\begin{tikzpicture}[xscale=0.9, yscale=0.85]
    % Axis
    \draw[thick, -Stealth] (0,0) -- (12,0) node[right, font=\small] {Values};
    \foreach \x in {1,3,5,7,9,11} \draw (\x, 0.1) -- (\x, -0.1) node[below, font=\scriptsize] {\x};

    % Box (Q1=4, Mediana=6, Q3=8)
    \draw[thick, fill=blue!15] (4,1) rectangle (8,3);
    \draw[very thick, red!80!black] (6,1) -- (6,3) node[above, font=\scriptsize\bfseries] {Median ($Q_2$)};

    % Whiskers
    \draw[thick] (4,2) -- (2,2);
    \draw[thick] (2,1.5) -- (2,2.5) node[above, font=\scriptsize] {Min Whisker};
    \draw[thick] (8,2) -- (10,2);
    \draw[thick] (10,1.5) -- (10,2.5) node[above, font=\scriptsize] {Max Whisker};

    % Outliers
    \fill[red] (11.2, 2) circle (3pt) node[above, font=\scriptsize] {Outlier};
    \fill[red] (0.8, 2) circle (3pt) node[above, font=\scriptsize] {Outlier};

    % Annotations
    \draw[Stealth-Stealth, thick] (4,0.5) -- node[below, font=\scriptsize] {$\text{IQR} = Q_3 - Q_1$} (8,0.5);
    \node[font=\scriptsize] at (4,3.3) {$Q_1$ (25\%)};
    \node[font=\scriptsize] at (8,3.3) {$Q_3$ (75\%)};
\end{tikzpicture}
\caption{Geometric anatomy of a Box Plot with identification of outliers.}
\label{fig:boxplot-anatomy}
\end{figure}

The limits of the whiskers are conventionally calculated by extending up to the most extreme observed value within the interval $[Q_1 - 1.5 \times \text{IQR}, \; Q_3 + 1.5 \times \text{IQR}]$. The observations that fall outside such thresholds are visually labeled as \textbf{outliers}.

A key extension for bivariate and multivariate exploration is given by \textbf{Conditional Box Plots}: the distribution of a continuous numerical attribute is stratified and displayed by placing box plots side by side for each category of a discrete attribute (typically the target class or a natural clustering). Figure~\ref{fig:conditional-boxplot} illustrates the distribution of \textit{Petal Length} stratified across the three flower species of the Iris dataset, highlighting that \textit{Setosa} is distinctly separable with very low dispersion, whereas \textit{Versicolor} and \textit{Virginica} exhibit a slight overlap in their interquartile ranges.

\begin{figure}[htbp]
\centering
\begin{tikzpicture}[xscale=0.95, yscale=0.85]
    % Axis
    \draw[thick, -Stealth] (0,0) -- (8.5,0) node[right, font=\small] {Petal Length (cm)};
    \foreach \x in {1,2,3,4,5,6,7,8} \draw (\x, 0.1) -- (\x, -0.1) node[below, font=\scriptsize] {\x};

    % Class 1: Setosa (around 1.5, narrow)
    \draw[thick, fill=blue!20] (1.3, 2.7) rectangle (1.7, 3.3);
    \draw[very thick, blue!90!black] (1.5, 2.7) -- (1.5, 3.3);
    \draw[thick] (1.3, 3.0) -- (1.0, 3.0); \draw[thick] (1.0, 2.8) -- (1.0, 3.2);
    \draw[thick] (1.7, 3.0) -- (1.9, 3.0); \draw[thick] (1.9, 2.8) -- (1.9, 3.2);
    \node[left, font=\small\bfseries, blue!80!black] at (0.8, 3.0) {Setosa};

    % Class 2: Versicolor (around 4.3, medium)
    \draw[thick, fill=orange!25] (4.0, 1.7) rectangle (4.7, 2.3);
    \draw[very thick, orange!90!black] (4.35, 1.7) -- (4.35, 2.3);
    \draw[thick] (4.0, 2.0) -- (3.3, 2.0); \draw[thick] (3.3, 1.8) -- (3.3, 2.2);
    \draw[thick] (4.7, 2.0) -- (5.1, 2.0); \draw[thick] (5.1, 1.8) -- (5.1, 2.2);
    \node[left, font=\small\bfseries, orange!80!black] at (0.8, 2.0) {Versicolor};

    % Class 3: Virginica (around 5.6, large)
    \draw[thick, fill=green!25] (5.1, 0.7) rectangle (6.0, 1.3);
    \draw[very thick, green!80!black] (5.55, 0.7) -- (5.55, 1.3);
    \draw[thick] (5.1, 1.0) -- (4.5, 1.0); \draw[thick] (4.5, 0.8) -- (4.5, 1.2);
    \draw[thick] (6.0, 1.0) -- (6.9, 1.0); \draw[thick] (6.9, 0.8) -- (6.9, 1.2);
    \node[left, font=\small\bfseries, green!60!black] at (0.8, 1.0) {Virginica};
\end{tikzpicture}
\caption{Example of Conditional Box Plot: distribution of the \textit{Petal Length} attribute conditioned on the flower species in the Iris dataset.}
\label{fig:conditional-boxplot}
\end{figure}

\begin{alertblock}{Critical Comparison: Box Plot vs Histogram}
Although the box plot effectively synthesizes position and dispersion, \textbf{it is not able to show the real shape of the distribution}. Two distributions having radically different densities (for example, a uniform unimodal distribution and a strongly bimodal distribution with an empty central valley) can produce identical values of $Q_1$, median, $Q_3$, Minimum, and Maximum, and therefore generate the exact same box plot. The histogram is essential to detect multimodality.
\end{alertblock}

\begin{figure}[htbp]
\centering
\begin{minipage}{0.46\textwidth}
\centering
\begin{tikzpicture}[scale=0.65]
    \draw[thick, -Stealth] (0,0) -- (5.5,0) node[right, font=\tiny] {$x$};
    \draw[thick, -Stealth] (0,0) -- (0,3.5) node[above, font=\tiny] {Frequency};
    % Uniform / unimodal shape
    \draw[fill=orange!30, draw=orange!80, thick] (0.5,0) rectangle (1.3, 1.0);
    \draw[fill=orange!30, draw=orange!80, thick] (1.3,0) rectangle (2.1, 2.2);
    \draw[fill=orange!30, draw=orange!80, thick] (2.1,0) rectangle (2.9, 3.0);
    \draw[fill=orange!30, draw=orange!80, thick] (2.9,0) rectangle (3.7, 2.2);
    \draw[fill=orange!30, draw=orange!80, thick] (3.7,0) rectangle (4.5, 1.0);
    \node[font=\scriptsize\bfseries] at (2.5, -0.6) {Unimodal Density};
\end{tikzpicture}
\end{minipage}\hfill
\begin{minipage}{0.46\textwidth}
\centering
\begin{tikzpicture}[scale=0.65]
    \draw[thick, -Stealth] (0,0) -- (5.5,0) node[right, font=\tiny] {$x$};
    \draw[thick, -Stealth] (0,0) -- (0,3.5) node[above, font=\tiny] {Frequency};
    % Bimodal shape with empty center
    \draw[fill=purple!30, draw=purple!80, thick] (0.5,0) rectangle (1.3, 2.8);
    \draw[fill=purple!30, draw=purple!80, thick] (1.3,0) rectangle (2.1, 1.2);
    \draw[fill=purple!30, draw=purple!80, thick] (2.1,0) rectangle (2.9, 0.3);
    \draw[fill=purple!30, draw=purple!80, thick] (2.9,0) rectangle (3.7, 1.2);
    \draw[fill=purple!30, draw=purple!80, thick] (3.7,0) rectangle (4.5, 2.8);
    \node[font=\scriptsize\bfseries] at (2.5, -0.6) {Bimodal Density};
\end{tikzpicture}
\end{minipage}

\vspace{0.3cm}
\begin{tikzpicture}[scale=0.75]
    \draw[thick, -Stealth] (0,0) -- (8.5,0) node[right, font=\scriptsize] {$x$};
    \foreach \x/\val in {0.8/Min, 2.4/$Q_1$, 4.0/Median, 5.6/$Q_3$, 7.2/Max} {
        \draw (\x, 0.1) -- (\x, -0.1) node[below, font=\tiny] {\val};
    }
    \draw[thick, fill=gray!25] (2.4,0.4) rectangle (5.6,1.2);
    \draw[very thick, red!80!black] (4.0,0.4) -- (4.0,1.2);
    \draw[thick] (2.4,0.8) -- (0.8,0.8); \draw[thick] (0.8,0.5) -- (0.8,1.1);
    \draw[thick] (5.6,0.8) -- (7.2,0.8); \draw[thick] (7.2,0.5) -- (7.2,1.1);
    \node[font=\scriptsize, red!80!black] at (4.0, 1.55) {Identical Box Plot for both distributions!};
\end{tikzpicture}
\caption{Critical comparison: two distributions with diametrically opposite behaviors (unimodal vs bimodal with empty center) generate the same five-number summary and the same Box Plot.}
\label{fig:boxplot-vs-hist}
\end{figure}

\subsection{Bar Chart for Categorical Attributes}
The \textbf{Bar Chart} visually represents the frequency distribution of a qualitative or categorical attribute. Each discrete category is represented by a vertical (or horizontal) bar, whose height corresponds to the absolute count or relative frequency. Unlike histograms, bars are separated by empty spaces to highlight that the modalities do not constitute a numerical continuum.

\begin{figure}[htbp]
\centering
\begin{tikzpicture}[xscale=1.4, yscale=0.07]
    \draw[thick, -Stealth] (0,0) -- (4.5,0) node[right, font=\small] {Marital Status};
    \draw[thick, -Stealth] (0,0) -- (0,55) node[above, font=\small] {Frequency};
    \foreach \y in {10, 20, 30, 40, 50} \draw (0.08, \y) -- (-0.08, \y) node[left, font=\scriptsize] {\y};

    \draw[fill=blue!35, draw=blue!80, thick] (0.6,0) rectangle (1.4,42);
    \node[above, font=\scriptsize\bfseries] at (1.0, 42) {42};
    \node[below, font=\small] at (1.0, -2) {Single};

    \draw[fill=blue!35, draw=blue!80, thick] (1.8,0) rectangle (2.6,33);
    \node[above, font=\scriptsize\bfseries] at (2.2, 33) {33};
    \node[below, font=\small] at (2.2, -2) {Married};

    \draw[fill=blue!35, draw=blue!80, thick] (3.0,0) rectangle (3.8,18);
    \node[above, font=\scriptsize\bfseries] at (3.4, 18) {18};
    \node[below, font=\small] at (3.4, -2) {Divorced};
\end{tikzpicture}
\caption{Example of a Bar Chart for a nominal categorical attribute (Marital Status).}
\label{fig:barchart-example}
\end{figure}

\subsection{Histograms for Numerical Attributes and Choice of Bins}
A \textbf{histogram} shows the empirical frequency distribution of a continuous numerical attribute. The domain of variability of the quantity is partitioned (\textit{discretized}) into a predetermined number $k$ of contiguous intervals called \textit{bins} (or class widths). For each bin, the height of the overlying bar is proportional to the number of sample observations that fall within it.

\begin{figure}[htbp]
\centering
\begin{tikzpicture}[xscale=0.9, yscale=0.08]
    \draw[thick, -Stealth] (0,0) -- (11.5,0) node[right, font=\small] {Age};
    \draw[thick, -Stealth] (0,0) -- (0,50) node[above, font=\small] {Frequency};
    \foreach \y in {10, 20, 30, 40, 50} \draw (0.1, \y) -- (-0.1, \y) node[left, font=\scriptsize] {\y};

    % 6 contiguous bins: 20-30, 30-40, 40-50, 50-60, 60-70, 70-80
    \draw[fill=teal!30, draw=teal!80, thick] (1,0) rectangle (2.5, 14);
    \draw[fill=teal!30, draw=teal!80, thick] (2.5,0) rectangle (4, 44);
    \draw[fill=teal!30, draw=teal!80, thick] (4,0) rectangle (5.5, 32);
    \draw[fill=teal!30, draw=teal!80, thick] (5.5,0) rectangle (7, 20);
    \draw[fill=teal!30, draw=teal!80, thick] (7,0) rectangle (8.5, 11);
    \draw[fill=teal!30, draw=teal!80, thick] (8.5,0) rectangle (10, 5);

    \foreach \x/\label in {1/20, 2.5/30, 4/40, 5.5/50, 7/60, 8.5/70, 10/80} {
        \draw (\x, 1.5) -- (\x, -1.5) node[below, font=\scriptsize] {\label};
    }
\end{tikzpicture}
\caption{Example of a Histogram for a continuous numerical attribute with contiguous bins of constant width.}
\label{fig:histogram-example}
\end{figure}

An incorrect choice of the number of bins profoundly alters the perception of the data:
\begin{itemize}
    \item An insufficient number of bins causes \textit{oversmoothing}, leveling the distribution and masking skewnesses and multimodality.
    \item An excessive number of bins causes \textit{undersmoothing}, fragmenting the sample into discontinuous frequencies and emphasizing sampling noise.
\end{itemize}

For approximately normal unimodal distributions and moderate sample sizes $n$, the standard estimate for the optimal number of intervals is given by \textbf{Sturges' rule}:
\begin{equation}
k = 1 + \lceil \log_2(n) \rceil = 1 + \lceil 3.322 \log_{10}(n) \rceil
\end{equation}
where $n$ represents the sample cardinality. This heuristic explicitly assumes a Gaussian-like distribution. When the sample size grows considerably or the underlying distribution exhibits multiple modes or heavy skewness, Sturges' rule tends to generate an insufficient number of bins; in such cases, it is crucial to test alternative values of $k$ experimentally (e.g., comparing low-, medium-, and high-resolution partitionings such as 5, 17, or 50 bins) to assess the robustness of observed patterns.

\section{Multivariate Visualization Techniques}

\subsection{Scatter Plot and Multi-Dimensional Visualization}
The \textbf{Scatter Plot} treats pairs of numerical attributes as Cartesian coordinates $(x_i, y_i)$ in the plane. It is the main visual tool for:
\begin{itemize}
    \item Identifying linear correlations or non-linear functional trends.
    \item Detecting natural clusters of points (\textit{clusters}).
    \item Finding multivariate outliers not visible in univariate projections.
\end{itemize}

\begin{example}[Analysis of the Anderson/Fisher Iris Dataset]
The famous \textit{Iris} dataset (150 samples of flowers belonging to three species: \textit{Setosa}, \textit{Versicolor}, \textit{Virginica}) describes 4 numerical attributes: sepal length and width, petal length and width.
The bivariate scatter plot shows that the attributes related to the petal (\textit{Petal Length} vs \textit{Petal Width}) guarantee a clear separation among the species, also allowing the isolation of potential anomalies.
\end{example}

\begin{figure}[htbp]
\centering
\begin{tikzpicture}[scale=0.95]
    % Axes
    \draw[thick, -Stealth] (0,0) -- (7.5,0) node[right, font=\small] {Petal Length (cm)};
    \draw[thick, -Stealth] (0,0) -- (0,5.5) node[above, font=\small] {Petal Width (cm)};
    \foreach \x in {1,2,3,4,5,6,7} \draw (\x, 0.1) -- (\x, -0.1) node[below, font=\scriptsize] {\x};
    \foreach \y in {1,2,3,4,5} \draw (0.1, \y) -- (-0.1, \y) node[left, font=\scriptsize] {\y};

    % Setosa points (bottom-left, blue circles)
    \foreach \p in {(1.2, 0.4), (1.4, 0.3), (1.5, 0.5), (1.6, 0.4), (1.3, 0.6), (1.7, 0.5), (1.4, 0.7)} {
        \fill[blue!75] \p circle (2.5pt);
    }
    \node[blue!80!black, font=\scriptsize\bfseries] at (1.5, 1.0) {Setosa};

    % Outlier for Setosa
    \draw[red!80!black, thick] (2.5, 1.2) circle (6pt);
    \fill[blue!75] (2.5, 1.2) circle (2.5pt);
    \node[red!80!black, font=\tiny, align=center] at (2.7, 1.7) {Local Outlier\\(for Setosa)};

    % Versicolor points (center, orange squares)
    \foreach \p in {(4.0, 2.2), (4.3, 2.4), (4.5, 2.8), (4.2, 2.6), (4.7, 2.5), (4.4, 2.3)} {
        \fill[orange!85!black] \p rectangle ++(4.5pt, 4.5pt);
    }
    \node[orange!80!black, font=\scriptsize\bfseries] at (4.5, 1.8) {Versicolor};

    % Virginica points (top-right, green triangles)
    \foreach \p in {(5.5, 3.8), (5.8, 4.2), (6.1, 4.5), (5.9, 4.0), (6.4, 4.4), (6.2, 4.1)} {
        \fill[green!50!black] \p ++(-2.5pt, -2.5pt) -- ++(5pt, 0pt) -- ++(-2.5pt, 5pt) -- cycle;
    }
    \node[green!50!black, font=\scriptsize\bfseries] at (6.0, 3.4) {Virginica};

    % Global Outlier
    \draw[red, very thick] (6.8, 0.9) circle (6pt);
    \fill[red] (6.8, 0.9) circle (3pt);
    \node[red!80!black, font=\tiny, align=center] at (6.8, 0.4) {\textbf{Global Outlier}};
\end{tikzpicture}
\caption{Bivariate Scatter Plot with graphical class encoding and visual detection of local and global outliers.}
\label{fig:scatterplot-outlier}
\end{figure}

To simultaneously explore all pairs of numerical attributes, a \textbf{Scatter Plot Matrix} (SPLOM) is used, which is an $n \times n$ grid where each off-diagonal cell $(i, j)$ displays the scatter plot between attribute $i$ and attribute $j$, while the main diagonal hosts the histograms or univariate density estimations.

\begin{figure}[htbp]
\centering
\begin{tikzpicture}[scale=0.92]
    % Cell dimensions for 3x3 SPLOM: Sepal Length (0), Petal Length (1), Petal Width (2)
    \def\cw{2.5}
    \def\ch{2.5}
    \def\gap{0.35}

    % Draw 3x3 Grid
    \foreach \col in {0, 1, 2} {
        \foreach \row in {0, 1, 2} {
            \draw[gray!40, thick, fill=white] (\col*\cw+\col*\gap, \row*\ch+\row*\gap) rectangle ++(\cw, \ch);
        }
    }

    % --- Diagonal: Feature Names and Density Approximations ---
    % Cell (0, 2) Top-Left: Sepal Length
    \node[font=\footnotesize\bfseries, align=center] at (0*\cw+0*\gap + 0.5*\cw, 2*\ch+2*\gap + 1.8) {Sepal\\Length};
    \draw[thick, fill=blue!12] (0*\cw+0*\gap + 0.3, 2*\ch+2*\gap + 0.3) to[bend left=35] ++(0.95, 0.95) to[bend left=35] ++(0.95, -0.95) -- cycle;

    % Cell (1, 1) Center: Petal Length
    \node[font=\footnotesize\bfseries, align=center] at (1*\cw+1*\gap + 0.5*\cw, 1*\ch+1*\gap + 1.8) {Petal\\Length};
    \draw[thick, fill=blue!12] (1*\cw+1*\gap + 0.3, 1*\ch+1*\gap + 0.3) to[bend left=50] ++(0.65, 0.9) to[bend left=30] ++(0.6, -0.55) to[bend left=35] ++(0.65, -0.35) -- cycle;

    % Cell (2, 0) Bottom-Right: Petal Width
    \node[font=\footnotesize\bfseries, align=center] at (2*\cw+2*\gap + 0.5*\cw, 0*\ch+0*\gap + 1.8) {Petal\\Width};
    \draw[thick, fill=blue!12] (2*\cw+2*\gap + 0.3, 0*\ch+0*\gap + 0.3) to[bend left=50] ++(0.6, 0.9) to[bend left=35] ++(0.65, -0.5) to[bend left=40] ++(0.65, -0.4) -- cycle;

    % --- Off-Diagonal Cells: Scatter of all 3 Iris species ---
    % Col 0, Row 1: X=SL, Y=PL
    \foreach \x/\y in {0.5/0.4, 0.8/0.5, 1.1/0.4, 0.7/0.6, 1.3/0.5} \fill[blue!75] (0*\cw+0*\gap + \x, 1*\ch+1*\gap + \y) circle (2pt);
    \foreach \x/\y in {1.2/1.3, 1.4/1.5, 1.6/1.4, 1.8/1.6, 1.5/1.7} \node[rectangle, fill=orange!85!black, inner sep=1.5pt] at (0*\cw+0*\gap + \x, 1*\ch+1*\gap + \y) {};
    \foreach \x/\y in {1.7/1.9, 1.9/2.2, 2.1/2.0, 2.3/2.3, 2.0/2.1} \node[regular polygon, regular polygon sides=3, fill=green!50!black, inner sep=1.3pt] at (0*\cw+0*\gap + \x, 1*\ch+1*\gap + \y) {};

    % Col 1, Row 2: X=PL, Y=SL
    \foreach \y/\x in {0.5/0.4, 0.8/0.5, 1.1/0.4, 0.7/0.6, 1.3/0.5} \fill[blue!75] (1*\cw+1*\gap + \x, 2*\ch+2*\gap + \y) circle (2pt);
    \foreach \y/\x in {1.2/1.3, 1.4/1.5, 1.6/1.4, 1.8/1.6, 1.5/1.7} \node[rectangle, fill=orange!85!black, inner sep=1.5pt] at (1*\cw+1*\gap + \x, 2*\ch+2*\gap + \y) {};
    \foreach \y/\x in {1.7/1.9, 1.9/2.2, 2.1/2.0, 2.3/2.3, 2.0/2.1} \node[regular polygon, regular polygon sides=3, fill=green!50!black, inner sep=1.3pt] at (1*\cw+1*\gap + \x, 2*\ch+2*\gap + \y) {};

    % Col 0, Row 0: X=SL, Y=PW
    \foreach \x/\y in {0.5/0.3, 0.8/0.4, 1.1/0.3, 0.7/0.5, 1.3/0.4} \fill[blue!75] (0*\cw+0*\gap + \x, 0*\ch+0*\gap + \y) circle (2pt);
    \foreach \x/\y in {1.2/1.2, 1.4/1.4, 1.6/1.3, 1.8/1.5, 1.5/1.4} \node[rectangle, fill=orange!85!black, inner sep=1.5pt] at (0*\cw+0*\gap + \x, 0*\ch+0*\gap + \y) {};
    \foreach \x/\y in {1.7/1.9, 1.9/2.2, 2.1/2.0, 2.3/2.3, 2.0/2.1} \node[regular polygon, regular polygon sides=3, fill=green!50!black, inner sep=1.3pt] at (0*\cw+0*\gap + \x, 0*\ch+0*\gap + \y) {};

    % Col 2, Row 2: X=PW, Y=SL
    \foreach \y/\x in {0.5/0.3, 0.8/0.4, 1.1/0.3, 0.7/0.5, 1.3/0.4} \fill[blue!75] (2*\cw+2*\gap + \x, 2*\ch+2*\gap + \y) circle (2pt);
    \foreach \y/\x in {1.2/1.2, 1.4/1.4, 1.6/1.3, 1.8/1.5, 1.5/1.4} \node[rectangle, fill=orange!85!black, inner sep=1.5pt] at (2*\cw+2*\gap + \x, 2*\ch+2*\gap + \y) {};
    \foreach \y/\x in {1.7/1.9, 1.9/2.2, 2.1/2.0, 2.3/2.3, 2.0/2.1} \node[regular polygon, regular polygon sides=3, fill=green!50!black, inner sep=1.3pt] at (2*\cw+2*\gap + \x, 2*\ch+2*\gap + \y) {};

    % Col 1, Row 0: X=PL, Y=PW (separazione ottimale)
    \foreach \x/\y in {0.4/0.3, 0.5/0.4, 0.6/0.3, 0.4/0.5, 0.7/0.4} \fill[blue!75] (1*\cw+1*\gap + \x, 0*\ch+0*\gap + \y) circle (2pt);
    \foreach \x/\y in {1.3/1.1, 1.5/1.3, 1.4/1.2, 1.6/1.4, 1.7/1.3} \node[rectangle, fill=orange!85!black, inner sep=1.5pt] at (1*\cw+1*\gap + \x, 0*\ch+0*\gap + \y) {};
    \foreach \x/\y in {1.9/1.8, 2.1/2.1, 2.0/1.9, 2.3/2.3, 2.2/2.1} \node[regular polygon, regular polygon sides=3, fill=green!50!black, inner sep=1.3pt] at (1*\cw+1*\gap + \x, 0*\ch+0*\gap + \y) {};

    % Col 2, Row 1: X=PW, Y=PL
    \foreach \y/\x in {0.4/0.3, 0.5/0.4, 0.6/0.3, 0.4/0.5, 0.7/0.4} \fill[blue!75] (2*\cw+2*\gap + \x, 1*\ch+1*\gap + \y) circle (2pt);
    \foreach \y/\x in {1.3/1.1, 1.5/1.3, 1.4/1.2, 1.6/1.4, 1.7/1.3} \node[rectangle, fill=orange!85!black, inner sep=1.5pt] at (2*\cw+2*\gap + \x, 1*\ch+1*\gap + \y) {};
    \foreach \y/\x in {1.9/1.8, 2.1/2.1, 2.0/1.9, 2.3/2.3, 2.2/2.1} \node[regular polygon, regular polygon sides=3, fill=green!50!black, inner sep=1.3pt] at (2*\cw+2*\gap + \x, 1*\ch+1*\gap + \y) {};

    % --- Centered Legend below Matrix ---
    \node[anchor=north] at (1*\cw+1*\gap + 0.5*\cw, -0.35) {
        \begin{tikzpicture}[font=\footnotesize]
            \fill[blue!75] (0,0) circle (2.5pt);
            \node[right] at (0.15,0) {Iris-setosa};
            \node[rectangle, fill=orange!85!black, inner sep=2pt] at (2.4,0) {};
            \node[right] at (2.6,0) {Iris-versicolor};
            \node[regular polygon, regular polygon sides=3, fill=green!50!black, inner sep=1.5pt] at (5.2,0) {};
            \node[right] at (5.45,0) {Iris-virginica};
        \end{tikzpicture}
    };
\end{tikzpicture}
\caption{Scatter Plot Matrix (SPLOM) of the Iris dataset with the pairwise comparison for all three species and the univariate density estimations along the diagonal.}
\label{fig:splom-example}
\end{figure}

\subsection{Visualization as a Test and Exploratory Heuristic}
Visual inspection serves as an irreplaceable preliminary test on data structure, whose sound interpretation requires precise epistemological awareness:
\begin{itemize}
    \item \textbf{Presence of visual evidence:} When visualizations expose distinct patterns, separated clusters, or isolated exceptions (\textit{outliers}), there is empirical certainty that genuine structural properties or real phenomena exist in the dataset.
    \item \textbf{Absence of visual evidence does not imply absence of structure:} Conversely, when conventional visualizations display no obvious arrangement or appear purely random, one \textbf{cannot} infer that the data is devoid of relationships. Complex multivariate dependencies, subtle non-linear associations, or high-order interactions involving four or more attributes simultaneously elude two-dimensional projections, requiring higher-dimensional projections or analytical modeling techniques.
\end{itemize}

\subsection{High-Dimensionality Visualization}

\subsubsection{Parallel Coordinates}
When the number of dimensions exceeds 3, the conventional Cartesian representation becomes inapplicable. \textbf{Parallel Coordinates} replace the orthogonal axes with a set of equidistant vertical parallel axes, each corresponding to an attribute.

\begin{figure}[htbp]
\centering
\begin{tikzpicture}
    % Definizione coordinate orizzontali ben spaziate per evitare qualsiasi sovrapposizione
    \def\xSL{0.0}
    \def\xSW{3.3}
    \def\xPL{6.6}
    \def\xPW{9.9}
    \def\h{5.0}

    % 4 vertical axes
    \draw[thick] (\xSL, 0) -- (\xSL, \h) node[above=8pt, font=\small\bfseries, align=center] {Sepal\\Length};
    \draw[thick] (\xSW, 0) -- (\xSW, \h) node[above=8pt, font=\small\bfseries, align=center] {Sepal\\Width};
    \draw[thick] (\xPL, 0) -- (\xPL, \h) node[above=8pt, font=\small\bfseries, align=center] {Petal\\Length};
    \draw[thick] (\xPW, 0) -- (\xPW, \h) node[above=8pt, font=\small\bfseries, align=center] {Petal\\Width};

    % Ticks and normalized dimensions (0 - 1)
    \foreach \x in {\xSL, \xSW, \xPL, \xPW} {
        \foreach \y in {0.0, 1.25, 2.5, 3.75, 5.0} {
            \draw (\x-0.08, \y) -- (\x+0.08, \y);
        }
    }
    \node[left, font=\tiny, gray] at (\xSL, 0.0) {min};
    \node[left, font=\tiny, gray] at (\xSL, \h) {max};

    % --- Species 1: Setosa (Blue) ---
    % Characteristics: Sepal L medium-low, Sepal W high, Petal L minimum, Petal W minimum
    \draw[blue!80, thick] (\xSL, 1.4) -- (\xSW, 3.9) -- (\xPL, 0.8) -- (\xPW, 0.6);
    \draw[blue!80, thick] (\xSL, 1.8) -- (\xSW, 4.3) -- (\xPL, 1.0) -- (\xPW, 0.8);
    \foreach \x/\y in {\xSL/1.4, \xSW/3.9, \xPL/0.8, \xPW/0.6, \xSL/1.8, \xSW/4.3, \xPL/1.0, \xPW/0.8} \fill[blue!80] (\x, \y) circle (2.5pt);

    % --- Species 2: Versicolor (Orange) ---
    % Characteristics: Sepal L medium, Sepal W medium-low, Petal L intermediate, Petal W intermediate
    \draw[orange!90!black, thick] (\xSL, 2.7) -- (\xSW, 1.9) -- (\xPL, 2.8) -- (\xPW, 2.5);
    \draw[orange!90!black, thick] (\xSL, 3.0) -- (\xSW, 2.3) -- (\xPL, 3.1) -- (\xPW, 2.8);
    \foreach \x/\y in {\xSL/2.7, \xSW/1.9, \xPL/2.8, \xPW/2.5, \xSL/3.0, \xSW/2.3, \xPL/3.1, \xPW/2.8} \node[rectangle, fill=orange!90!black, inner sep=2pt] at (\x, \y) {};

    % --- Species 3: Virginica (Green) ---
    % Characteristics: Sepal L medium-high, Sepal W medium, Petal L high, Petal W high
    \draw[green!60!black, thick] (\xSL, 4.2) -- (\xSW, 2.6) -- (\xPL, 4.4) -- (\xPW, 4.1);
    \draw[green!60!black, thick] (\xSL, 4.6) -- (\xSW, 2.9) -- (\xPL, 4.8) -- (\xPW, 4.5);
    \foreach \x/\y in {\xSL/4.2, \xSW/2.6, \xPL/4.4, \xPW/4.1, \xSL/4.6, \xSW/2.9, \xPL/4.8, \xPW/4.5} \node[regular polygon, regular polygon sides=3, fill=green!60!black, inner sep=1.8pt] at (\x, \y) {};

    % --- Well-spaced legend below ---
    \node[anchor=north] at (5.0, -0.5) {
        \begin{tikzpicture}[font=\footnotesize\bfseries]
            \draw[blue!80, very thick] (0,0) -- (0.6,0); \fill[blue!80] (0.3,0) circle (2.5pt);
            \node[right, blue!80] at (0.65,0) {Iris-setosa};

            \draw[orange!90!black, very thick] (3.3,0) -- (3.9,0); \node[rectangle, fill=orange!90!black, inner sep=2pt] at (3.6,0) {};
            \node[right, orange!90!black] at (3.95,0) {Iris-versicolor};

            \draw[green!60!black, very thick] (6.8,0) -- (7.4,0); \node[regular polygon, regular polygon sides=3, fill=green!60!black, inner sep=1.8pt] at (7.1,0) {};
            \node[right, green!60!black] at (7.45,0) {Iris-virginica};
        \end{tikzpicture}
    };
\end{tikzpicture}
\caption{Representation of 4-dimensional data using Parallel Coordinates for all three species of the Iris dataset.}
\label{fig:parallel-coordinates}
\end{figure}

\begin{itemize}
    \item Each multidimensional object is represented by a \textbf{polygonal broken line} that connects the respective scaled values along the vertical axes.
    \item Objects belonging to the same class tend to form bundles of parallel lines or with coordinated trends.
    \item \textbf{Sensitivity to the ordering of the axes:} The visibility of clusters and correlations depends critically on the sequential arrangement of the axes. Two strongly correlated attributes show regular geometric patterns only if positioned as adjacent axes.
\end{itemize}

\subsubsection{Radar Plot and Star Plot}
In \textbf{Radar Plots}, the axes start from a common center and radiate out in a star shape at constant angular intervals. The value of each attribute defines a vertex along its own radius; joining the vertices yields a polygon for each instance.
\textbf{Star Plots} apply the same principle by drawing each object in a separate graph, facilitating comparison for individual profiles.

\begin{figure}[htbp]
\centering
\begin{tikzpicture}[scale=0.9]
    % Central origin
    \coordinate (O) at (0,0);
    % 4 axes at 90 degrees: North, East, South, West
    \draw[gray!60, -Stealth, thick] (0,0) -- (0,3.6) node[above=4pt, font=\scriptsize\bfseries] {Sepal Length};
    \draw[gray!60, -Stealth, thick] (0,0) -- (3.6,0) node[right=4pt, font=\scriptsize\bfseries] {Sepal Width};
    \draw[gray!60, -Stealth, thick] (0,0) -- (0,-3.6) node[below=4pt, font=\scriptsize\bfseries] {Petal Length};
    \draw[gray!60, -Stealth, thick] (0,0) -- (-3.6,0) node[left=4pt, font=\scriptsize\bfseries] {Petal Width};

    % Concentric reference circles
    \foreach \r in {1, 2, 3} {
        \draw[gray!25, dashed] (0,0) circle (\r cm);
    }

    % 1. Setosa (Blue): minimum petals, wide sepal width
    \coordinate (S1) at (0, 1.5);
    \coordinate (S2) at (2.9, 0);
    \coordinate (S3) at (0, -0.8);
    \coordinate (S4) at (-0.6, 0);
    \draw[blue!80, very thick, fill=blue!12, fill opacity=0.4] (S1) -- (S2) -- (S3) -- (S4) -- cycle;
    \foreach \p in {S1, S2, S3, S4} \fill[blue!80] (\p) circle (2.5pt);

    % 2. Versicolor (Orange): intermediate balanced values
    \coordinate (C1) at (0, 2.3);
    \coordinate (C2) at (1.4, 0);
    \coordinate (C3) at (0, -2.1);
    \coordinate (C4) at (-1.7, 0);
    \draw[orange!90!black, very thick, fill=orange!15, fill opacity=0.4] (C1) -- (C2) -- (C3) -- (C4) -- cycle;
    \foreach \p in {C1, C2, C3, C4} \fill[orange!90!black] (\p) ++(-2pt,-2pt) rectangle ++(4pt,4pt);

    % 3. Virginica (Green): maximum petals and high sepal length
    \coordinate (V1) at (0, 3.2);
    \coordinate (V2) at (1.8, 0);
    \coordinate (V3) at (0, -3.1);
    \coordinate (V4) at (-2.8, 0);
    \draw[green!60!black, very thick, fill=green!15, fill opacity=0.4] (V1) -- (V2) -- (V3) -- (V4) -- cycle;
    \foreach \p in {V1, V2, V3, V4} \fill[green!60!black] (\p) ++(-2.5pt, -2pt) -- ++(5pt, 0pt) -- ++(-2.5pt, 4.5pt) -- cycle;

    % Legend for all three species
    \node[anchor=north] at (0, -4.1) {
        \begin{tikzpicture}[font=\footnotesize\bfseries]
            \fill[blue!80] (0,0) circle (2.5pt);
            \node[right, blue!80] at (0.2,0) {Iris-setosa};

            \node[rectangle, fill=orange!90!black, inner sep=2pt] at (2.8,0) {};
            \node[right, orange!90!black] at (3.05,0) {Iris-versicolor};

            \node[regular polygon, regular polygon sides=3, fill=green!60!black, inner sep=1.8pt] at (5.9,0) {};
            \node[right, green!60!black] at (6.2,0) {Iris-virginica};
        \end{tikzpicture}
    };
\end{tikzpicture}
\caption{Example of a Radar Plot with visual comparison of multivariate profiles for all three Iris species.}
\label{fig:radar-plot-example}
\end{figure}

\section{Linear Correlation and Association Analysis}

Correlation analysis aims to identify, quantify, and interpret the strength and direction of the linear relationship between pairs of continuous numerical attributes.

\subsection{Motivations and Multicollinearity in Predictive Models}
Within the data understanding workflow, assessing correlation between variables serves several methodological goals:
\begin{itemize}
    \item \textbf{Identifying informative redundancies:} If two attributes exhibit nearly identical trajectories, one of them can be discarded without appreciable loss of descriptive power, reducing the dimensionality of the feature space.
    \item \textbf{Preparation for predictive models and clustering:} Many machine learning algorithms (such as linear regression, logistic regression, or distance-based clustering) assume independent features or benefit substantially from the removal of collinear attributes.
    \item \textbf{Detecting multicollinearity:} Occurs when two or more predictor variables (\textit{features}) in a regression model are strongly linearly related to one another. Consider the paradigmatic example of predicting house prices using simultaneously \textit{surface area in square metres} ($m^2$) and \textit{surface area in square feet} ($\text{ft}^2$). Both variables convey exactly the same information up to a constant conversion factor. Consequently, the regression model experiences severe numerical instability, struggling to disentangle the individual marginal contributions of each predictor and inflating estimator variance.
\end{itemize}

\subsection{From Variance to Sample Covariance}
While sample variance quantifies the dispersion of observations for a single variable around its own mean:
\begin{equation}
s_x^2 = \operatorname{Var}(X) = \frac{1}{n-1} \sum_{i=1}^n (x_i - \bar{x})^2
\end{equation}
in data mining applications interest is constantly focused on pairwise interactions between attributes:
\begin{itemize}
    \item Do customers with higher taxable incomes spend larger sums?
    \item Does a product's unit price lead to a contraction in sales volume?
    \item Do students who study more weekly hours systematically attain higher exam grades?
    \item Do two physical sensors arranged in parallel provide largely redundant measurement streams?
\end{itemize}

To quantify how two numerical quantities vary jointly, we introduce \textbf{sample covariance}:

\begin{definition}[Sample Covariance]
Given a sample of $n$ observation pairs $(x_i, y_i)$, the sample covariance between $X$ and $Y$ is defined as:
\begin{equation}
\operatorname{Cov}(X, Y) = \frac{1}{n-1} \sum_{i=1}^n (x_i - \bar{x})(y_i - \bar{y})
\end{equation}
where $\bar{x} = \frac{1}{n}\sum_{i=1}^n x_i$ and $\bar{y} = \frac{1}{n}\sum_{i=1}^n y_i$ represent the respective sample means.
\end{definition}

A fundamental theoretical property dictates that variance is simply the covariance of a variable with itself:
\begin{equation}
\operatorname{Cov}(X, X) = \frac{1}{n-1} \sum_{i=1}^n (x_i - \bar{x})^2 = \operatorname{Var}(X) = s_x^2
\end{equation}

\subsection{Sign Origin and Geometric Quadrant Interpretation}
The algebraic sign of each term in the summation stems from the product of paired deviations from their respective means:
\begin{itemize}
    \item \textbf{Positive concordant deviations:} $(x_i - \bar{x}) > 0$ and $(y_i - \bar{y}) > 0 \implies (+)\times(+) = \mathbf{+}$.
    \item \textbf{Negative concordant deviations:} $(x_i - \bar{x}) < 0$ and $(y_i - \bar{y}) < 0 \implies (-)\times(-) = \mathbf{+}$.
    \item \textbf{Discordant deviations:} $(+)\times(-) = \mathbf{-}$ or $(-)\times(+) = \mathbf{-}$.
\end{itemize}

Translating the Cartesian origin to the sample centroid $(\bar{x}, \bar{y})$, the two-dimensional space is divided into four quadrants (Figure~\ref{fig:covariance-quadrants}):
\begin{enumerate}
    \item If observations in quadrants \textbf{I and III} (concordant deviations) dominate, positive products prevail and the sample covariance is strictly positive ($\operatorname{Cov}(X, Y) > 0$).
    \item If quadrants \textbf{II and IV} (discordant deviations) prevail, negative products dominate and the covariance is strictly negative ($\operatorname{Cov}(X, Y) < 0$).
    \item If observations are uniformly or symmetrically scattered across all four quadrants, positive and negative terms cancel each other out, resulting in a covariance close to zero ($\operatorname{Cov}(X, Y) \approx 0$).
\end{enumerate}

\begin{figure}[htbp]
\centering
\begin{tikzpicture}[scale=0.95]
    % Central cross at mean (x_bar, y_bar)
    \def\xb{3.2}
    \def\yb{2.8}

    % Shaded quadrants
    \fill[green!8] (\xb, \yb) rectangle (6.5, 5.3);
    \fill[green!8] (0.5, 0.5) rectangle (\xb, \yb);
    \fill[red!8] (0.5, \yb) rectangle (\xb, 5.3);
    \fill[red!8] (\xb, 0.5) rectangle (6.5, \yb);

    % Axes
    \draw[thick, -Stealth] (0.5, 0.5) -- (6.8, 0.5) node[right, font=\small] {$X$};
    \draw[thick, -Stealth] (0.5, 0.5) -- (0.5, 5.6) node[above, font=\small] {$Y$};

    % Mean dashed lines
    \draw[dashed, thick, gray!70] (\xb, 0.5) -- (\xb, 5.3);
    \draw[dashed, thick, gray!70] (0.5, \yb) -- (6.5, \yb);
    \node[below, font=\small\bfseries] at (\xb, 0.5) {$\bar{x}$};
    \node[left, font=\small\bfseries] at (0.5, \yb) {$\bar{y}$};

    % Quadrant Labels and signs
    \node[align=center, font=\scriptsize] at (5.0, 4.6) {\textbf{Quadrant I}\\$(x_i > \bar{x}, \; y_i > \bar{y})$\\$(+) \times (+) = \mathbf{+}$};
    \node[align=center, font=\scriptsize] at (1.8, 4.6) {\textbf{Quadrant II}\\$(x_i < \bar{x}, \; y_i > \bar{y})$\\$(-) \times (+) = \mathbf{-}$};
    \node[align=center, font=\scriptsize] at (1.8, 1.4) {\textbf{Quadrant III}\\$(x_i < \bar{x}, \; y_i < \bar{y})$\\$(-) \times (-) = \mathbf{+}$};
    \node[align=center, font=\scriptsize] at (5.0, 1.4) {\textbf{Quadrant IV}\\$(x_i > \bar{x}, \; y_i < \bar{y})$\\$(+) \times (-) = \mathbf{-}$};

    % Sample points illustrating positive covariance
    \foreach \p in {(1.2, 1.1), (1.6, 1.8), (2.1, 1.5), (2.5, 2.2), (2.8, 2.4),
                    (3.6, 3.2), (4.0, 3.8), (4.5, 3.5), (5.1, 4.2), (5.8, 4.7), (6.1, 4.9)} {
        \fill[blue!80!black] \p circle (2.5pt);
    }
    \fill[blue!80!black] (2.2, 3.3) circle (2.5pt);
    \fill[blue!80!black] (4.2, 2.0) circle (2.5pt);

    \node[font=\footnotesize\itshape, blue!80!black] at (3.5, -0.2) {Dominance in Quadrants I and III $\implies \operatorname{Cov}(X, Y) > 0$};
\end{tikzpicture}
\caption{Geometric interpretation of covariance across quadrants centered at $(\bar{x}, \bar{y})$. Observations in concordant quadrants I and III yield positive contributions, whereas discordant quadrants II and IV yield negative contributions.}
\label{fig:covariance-quadrants}
\end{figure}

\subsection{Limitation of Covariance: Scale and Unit Dependence}
Although covariance unambiguously reveals the direction of linear agreement, it suffers from a critical drawback: \textbf{it depends strictly on the measurement units of the variables}.
Consider the relationship between annual income and household expenditure:
\begin{itemize}
    \item If both income and expenditure are measured in euros, the unit of $\operatorname{Cov}(\text{Income}, \text{Expenditure})$ is $\text{euro}^2$.
    \item If income is converted into \textit{thousands of euros}, the numerical covariance drops by a factor of $1000$, even though the substantive economic relationship between both quantities remains strictly unchanged.
\end{itemize}
Consequently, covariance indicates whether the linear trend is positive or negative, but its absolute magnitude cannot be compared across different scales or attribute pairs without standardization.

\subsection{Pearson Linear Correlation Coefficient}
To obtain a unit-free, standardized metric, sample covariance is divided by the product of the respective sample standard deviations:

\begin{definition}[Pearson Correlation Coefficient]
Given a sample of $n$ pairs $(x_i, y_i)$, the sample linear correlation coefficient $r_{xy}$ is defined as:
\begin{equation}
r_{xy} = \frac{\operatorname{Cov}(X, Y)}{s_x s_y} = \frac{\sum_{i=1}^n (x_i - \bar{x})(y_i - \bar{y})}{\sqrt{\sum_{i=1}^n (x_i - \bar{x})^2 \sum_{i=1}^n (y_i - \bar{y})^2}}
\end{equation}
where $s_x, s_y$ are the sample standard deviations of $X$ and $Y$.
\end{definition}

Core properties of the Pearson correlation coefficient:
\begin{itemize}
    \item \textbf{Dimensionless and scale-invariant:} $r_{xy}$ is a pure number, invariant under any positive affine transformation ($X' = aX + b$, with $a > 0$).
    \item \textbf{Bounded range:} $-1 \le r_{xy} \le +1$.
    \item \textbf{Perfect collinearity ($|r_{xy}| = 1$):} Points lie precisely on a straight line; with positive slope if $r_{xy} = +1$, and negative slope if $r_{xy} = -1$.
    \item \textbf{Absence of linear association ($r_{xy} \approx 0$):} Signals little to no linear relationship.
    \item \textbf{Blindness to non-linear dependencies:} As illustrated in Figure~\ref{fig:pearson-cases}(d), zero correlation ($r_{xy} \approx 0$) does not imply statistical independence: deterministic quadratic or sinusoidal relationships may exhibit exactly zero Pearson correlation.
\end{itemize}

\begin{figure}[htbp]
\centering
\begin{minipage}{0.47\textwidth}
\centering
\begin{tikzpicture}[scale=0.55]
    \draw[thick, -Stealth] (0,0) -- (5,0) node[right, font=\tiny] {$X$};
    \draw[thick, -Stealth] (0,0) -- (0,5) node[above, font=\tiny] {$Y$};
    \foreach \p in {(0.5,0.7), (1,1.2), (1.5,1.4), (2,2.2), (2.5,2.6), (3,3.1), (3.5,3.7), (4,4.1), (4.5,4.6)} {
        \fill[blue!70] \p circle (2.5pt);
    }
    \draw[red, thick, dashed] (0.3,0.4) -- (4.7,4.8);
    \node[font=\scriptsize\bfseries] at (2.5,-0.8) {(a) $r_{xy} \approx +0.98$};
\end{tikzpicture}
\end{minipage}%
\hfill
\begin{minipage}{0.47\textwidth}
\centering
\begin{tikzpicture}[scale=0.55]
    \draw[thick, -Stealth] (0,0) -- (5,0) node[right, font=\tiny] {$X$};
    \draw[thick, -Stealth] (0,0) -- (0,5) node[above, font=\tiny] {$Y$};
    \foreach \p in {(0.5,4.5), (1,4.0), (1.5,3.6), (2,2.8), (2.5,2.5), (3,1.9), (3.5,1.4), (4,1.1), (4.5,0.5)} {
        \fill[blue!70] \p circle (2.5pt);
    }
    \draw[red, thick, dashed] (0.3,4.7) -- (4.7,0.3);
    \node[font=\scriptsize\bfseries] at (2.5,-0.8) {(b) $r_{xy} \approx -0.98$};
\end{tikzpicture}
\end{minipage}

\vspace{0.4cm}
\begin{minipage}{0.47\textwidth}
\centering
\begin{tikzpicture}[scale=0.55]
    \draw[thick, -Stealth] (0,0) -- (5,0) node[right, font=\tiny] {$X$};
    \draw[thick, -Stealth] (0,0) -- (0,5) node[above, font=\tiny] {$Y$};
    \foreach \p in {(1.2,3.5), (2.5,4.2), (3.8,3.2), (1.5,1.8), (2.8,2.2), (3.5,1.5), (2.0,2.8), (4.2,2.5), (1.0,2.5)} {
        \fill[blue!70] \p circle (2.5pt);
    }
    \node[font=\scriptsize\bfseries] at (2.5,-0.8) {(c) $r_{xy} \approx 0.0$ (No correlation)};
\end{tikzpicture}
\end{minipage}%
\hfill
\begin{minipage}{0.47\textwidth}
\centering
\begin{tikzpicture}[scale=0.55]
    \draw[thick, -Stealth] (0,0) -- (5,0) node[right, font=\tiny] {$X$};
    \draw[thick, -Stealth] (0,0) -- (0,5) node[above, font=\tiny] {$Y$};
    \foreach \p in {(0.5,4.0), (1.0,2.25), (1.5,1.0), (2.0,0.25), (2.5,0.05), (3.0,0.25), (3.5,1.0), (4.0,2.25), (4.5,4.0)} {
        \fill[red!80!black] \p circle (2.5pt);
    }
    \draw[red!80!black, thick] (0.5,4.0) parabola bend (2.5,0.05) (4.5,4.0);
    \node[font=\scriptsize\bfseries, align=center] at (2.5,-0.8) {(d) $r_{xy} \approx 0.0$ (Quadratic relation!)};
\end{tikzpicture}
\end{minipage}
\caption{Visual interpretation of the Pearson correlation coefficient: (a) positive correlation; (b) negative correlation; (c) absence of linear relation; (d) non-linear deterministic dependence missed by Pearson.}
\label{fig:pearson-cases}
\end{figure}

\subsection{Statistical Association vs Causation and Confounding Factors}
A fundamental tenet in data science is the epistemological boundary separating association from causation:
\begin{equation}
\textbf{Correlation} \;\neq\; \textbf{Causation}
\end{equation}
Consider an academic dataset where weekly study hours and examination marks exhibit a high positive correlation ($r_{xy} = +0.98$). Students with fewer study hours tend to score below the mean, while those studying longer hours score above it. Although this demonstrates an undeniably strong positive linear association, \textbf{it does not warrant the conclusion that 2 additional study hours will deterministically cause an increase in marks}. Multiple confounding variables (such as prior academic background, lecture attendance, or classroom attention) may influence both quantities simultaneously.

Linear correlation:
\begin{enumerate}
    \item Does not identify the causal direction ($X \to Y$ or $Y \to X$).
    \item Can be driven by a \textbf{confounding variable} (omitted variable bias).
\end{enumerate}

\begin{example}[The Confounding Variable: Ice Cream and Drowning]
In seasonal regional records, ice cream sales and drowning incidents show a strong positive correlation. Eating ice cream clearly does not cause drowning. The observed relationship is driven entirely by an omitted confounding variable: \textbf{high summer temperatures} (heatwaves), which simultaneously increase demand for chilled desserts and the number of people swimming, producing a spurious association.
\end{example}

\subsection{Spearman's Rank Correlation for Monotonic Relationships}
When the relationship between two variables is monotonic but non-linear, or in the presence of severe outliers that distort means and standard deviations, the \textbf{Spearman rank correlation coefficient} ($\rho_s$) is preferred.

Spearman does not evaluate original continuous values directly, but replaces each value with its \textbf{rank} in ascending order:

\begin{definition}[Spearman's Rank Correlation Coefficient]
Let $\operatorname{rank}(x_i)$ and $\operatorname{rank}(y_i)$ denote the ranks of $x_i$ and $y_i$. If there are no ties, the coefficient is given by:
\begin{equation}
\rho_s = 1 - \frac{6 \sum_{i=1}^n d_i^2}{n(n^2 - 1)}, \quad \text{with } d_i = \operatorname{rank}(x_i) - \operatorname{rank}(y_i)
\end{equation}
In the presence of tied ranks, $\rho_s$ is formally equivalent to Pearson's correlation computed directly on the rank vectors.
\end{definition}

Essential properties:
\begin{itemize}
    \item $\rho_s = +1$ when the rankings of both variables are in exactly the same order (\textbf{monotonically increasing} relationship).
    \item $\rho_s = -1$ when the rankings are in reverse order (\textbf{monotonically decreasing} relationship).
    \item Measures whether variables tend to move in the same or opposite direction regardless of linearity (e.g., for $Y = X^3$ or $Y = \log(X)$, $\rho_s = 1$).
    \item Highly robust against isolated extreme outliers.
\end{itemize}

\subsection{Correlation Matrix and Feature Selection}
In datasets with multiple continuous attributes, all pairwise bivariate correlations are systematically arranged in a symmetric $p \times p$ \textbf{correlation matrix}.
In the Iris dataset, for example, the matrix reveals an exceptionally strong correlation between \textit{Petal Length} and \textit{Petal Width} ($r \approx 0.96$), signaling feature redundancy and suggesting that dimensionality reduction or feature selection can discard one of the two without sacrificing model performance.

\section{Handling Outliers and Missing Values}

\subsection{Outliers: Noise vs Analytical Target}
An \textbf{outlier} is a data object whose characteristics deviate considerably from the vast majority of observations in the dataset:
\begin{itemize}
    \item \textbf{Outliers as noise:} Originating from typographic mistakes, sensor failures, or protocol errors. Such instances severely corrupt sample means, variances, and regression fits, and must be detected and removed (or imputed) during data cleaning.
    \item \textbf{Outliers as the primary analysis target:} Genuine unusual situations that constitute the central analytical goal. Paradigmatic examples include \textbf{credit card fraud detection}, \textbf{network intrusion detection}, and early diagnosis of \textbf{rare medical pathologies}.
\end{itemize}

Even when outliers represent valid measurements of rare events, it is often advisable to exclude them temporarily when computing summary metrics aimed at describing the typical population.

\subsection{Outlier Detection Methodologies}
Anomalies can be discovered across different dimensional settings:
\begin{itemize}
    \item \textbf{Single Attribute (Univariate Analysis):}
    \begin{itemize}
        \item \textit{Categorical attributes:} An outlier is a category that occurs with a frequency vastly lower than all other domain classes.
        \item \textit{Numerical attributes:} Visual and analytical inspection using Box Plots (points beyond $[Q_1 - 1.5 \times \text{IQR}, \; Q_3 + 1.5 \times \text{IQR}]$) and standardized $z$-scores ($z = \frac{x - \mu}{\sigma}$).
    \end{itemize}
    \item \textbf{Multiple Attributes (Multivariate Analysis):}
    \begin{itemize}
        \item \textit{2D and 3D Scatter Plots:} Reveal both global outliers (isolated with respect to the entire point cloud) and local outliers (deviating from their specific class cluster, as seen in Figure~\ref{fig:scatterplot-outlier}).
        \item \textit{PCA Projections:} Project high-dimensional data onto maximal-variance planes (Principal Component Analysis) to visually isolate anomalies.
        \item \textit{Clustering Algorithms:} In density-based or partition algorithms (such as DBSCAN or k-means with distance cutoffs), outliers naturally emerge as points that cannot be assigned to any major cluster or that form isolated micro-clusters.
    \end{itemize}
\end{itemize}

\subsection{Missing Values and Camouflage}
The absence of values for attributes across records arises from several root causes:
\begin{itemize}
    \item \textbf{Information not collected or withheld:} Survey respondents explicitly refuse to disclose sensitive personal data (e.g., income, weight, or age).
    \item \textbf{Logical non-applicability:} The attribute has no logical meaning for a subpopulation (e.g., annual employment income for school-age children, or driver's licenses for minors).
    \item \textbf{Sensor breakdown or network loss:} Physical hardware malfunctions, depleted sensor batteries, or packet dropouts.
\end{itemize}

\begin{alertblock}{Pitfall of Camouflaged Missing Values: Zeros and Default Sentinels}
Frequently, a missing value is not explicitly stored as \texttt{NULL} or \texttt{NaN}, but rather recorded using conventional default or sentinel numbers (such as \texttt{0}, \texttt{-1}, or \texttt{999}). A blocked or broken sensor, for instance, often transmits a constant stream of \texttt{0}s.
Unless intercepted during Data Understanding, analytics software treats these sentinel numbers as true physical measurements, severely biasing means, standard deviations, correlations, and classification models.
\end{alertblock}

\section{Operational Checklist for Data Understanding}
Upon concluding the exploration stage, the analyst must have methodically executed the following operational verifications:
\begin{enumerate}
    \item \textbf{Assess Overall Data Quality:} Audit syntactic accuracy (out-of-domain entries) and semantic consistency, check completeness, evaluate timeliness, deduplicate records, and quantify class imbalance.
    \item \textbf{Detect and Analyze Outliers:} Identify univariate and multivariate outliers using visualization tools (box plots, scatter plots, PCA, parallel coordinates) and determine appropriate handling policies (noise exclusion vs target investigation).
    \item \textbf{Uncover and Examine Missing Values:} Quantify null frequencies and expose camouflaged missing values hidden behind sentinel zeros or default settings.
    \item \textbf{Discover and Validate Correlations and Dependencies:} Compute linear (Pearson) and monotonic rank (Spearman) correlation matrices to confirm domain hypotheses, discover unexpected linkages, and prevent multicollinearity and feature redundancy.
    \item \textbf{Verify Domain-Specific Assumptions:} Test distributional premises required by downstream models or heuristics (e.g., verifying whether an attribute follows an approximately normal distribution before applying Sturges' rule or parametric tests).
    \item \textbf{Compare Statistics with Expected Behavior:} Ensure that empirical means, medians, skewness, and spreads align consistently with domain expertise and business context expectations.
\end{enumerate}
